\section{The difficulty on the terminal chain}\label{sec:difficulty}

We now construct an integer that is nonnegative at the crepant terminal
junctions and nonincreasing along the moves between them.  The local
subtraction of surface contributions and the correction by cycle ranks on
boundary normalizations adapt the projective construction in
\cite[Definition~4.2, Lemma~4.3, and Section~7]{ProjectiveLCFourfolds}.
We give the required finiteness and comparison arguments in the analytic
setting.

Use the terminal chain of Proposition~\ref{prop:terminal-junctions}, after
the truncation in Proposition~\ref{prop:surface-finiteness}.  Its finite set
of prime labels is denoted by \(\mathcal H\).  On a current model its
boundary is
\[
 \Theta=\sum_{h\in\mathcal H}b_hS_h,\qquad 0\leq b_h<1.
\]
The labels persist by strict transform, including those of coefficient
zero, and the coefficients are nonincreasing.  The fixed coefficients have
already become constant.  Every varying coefficient is positive; at a
junction \(h_i:Y_i\to X_i\), its label is exceptional for \(h_i\) and has
image of dimension at most one.  Each model of the chain is strong and
compact K\"ahler, the generalized pair is terminal, and the underlying
fourfold is ordinarily terminal.

Choose a positive integer \(N\) that clears the fixed label coefficients,
the coefficients of the initial boundary of the particular strong
generalized klt sequence under study, and its fixed nef datum: on the chosen
higher carrier, \(N\M\) is represented by an integral Cartier divisor.
These are finite rational data.  We impose no denominator condition on a
varying coefficient.  For \(u\in\R\), put
\begin{equation}\label{eq:difficulty-weights}
 \omega(u)=\left\lceil N\max\{u,0\}\right\rceil,
 \qquad
 w(E)=\omega\bigl(2-a(E;Y,\Theta+\M)\bigr),
\end{equation}
where \(E\) is a global divisorial place and the current pair is understood
in \(w(E)\).  Thus \(w(E)>0\) exactly when the discrepancy is less than
two; discrepancy equal to two has weight zero.

For a fixed label, \(\omega(b_h)=Nb_h\), so its contribution incurs no
rounding loss in the local estimate below.  The conditions on the initial
boundary and nef carrier will give the lattice at the positive endpoint in
Lemma~\ref{lem:positive-surface}.  The earlier surface lattice has already
served to choose the tail; \(N\) need not clear every element of
\(\mathcal A_{\mathrm{surf}}\).

\subsection{Echoes and finite local corrections}

The recurring places over a generic codimension-two boundary center are
the echoes of \cite[Example~1.4, Lemma~1.5, and Definition~2.3]{AlexeevHaconKawamata2007}.
The full chain matters: \cite[Example~2.1]{Fujino2005Addendum} shows why
removing only the first blowup does not make a boundary difficulty finite.

First consider an irreducible compact analytic surface \(V\subset S_h\)
whose general point lies in a smooth open of \(Y\) where \(S_h\) is smooth,
no other label occurs, and \(\M\) descends.  The following sequence of
global places records the predictable contribution above such a surface.
Blow up the coherent ideal of \(V\) globally and normalize.  The global
exceptional divisor has a unique irreducible component dominating the dense
smooth open of \(V\); call its place \(E_{V,h,1}\).  On this normalized
blowup, its analytic intersection with the strict transform of \(S_h\)
has a unique irreducible component dominating that same open of \(V\).
Blow up the coherent ideal of this global component, normalize, and take
the unique new exceptional prime over that open.  Repeat this operation on
these normalized blowups, denoting the successive places by
\(E_{V,h,j}\), \(j\geq1\).  The models are smooth at the general points
of the indicated intersections, so each next component exists and is
unique there.  When another higher model is needed for comparison, a common
proper model identifies these same global places by strict transform.  We
call them the \emph{echoes} of \(S_h\) above \(V\).

The first exceptional prime has crepant boundary coefficient
\(-(1-b_h)\).  At stage \(j\) this coefficient is \(-j(1-b_h)\): blowing
up its intersection with the strict label adds one more \(-(1-b_h)\).
Consequently
\begin{equation}\label{eq:echo-discrepancy}
 a(E_{V,h,j};Y,\Theta+\M)=1+j(1-b_h).
\end{equation}
The contributing echoes are precisely those with \(j(1-b_h)<1\).  In
particular there are only finitely many contributing indices for each
\(b_h<1\).  Define the \emph{default contribution} of label \(h\) by the finite sum
\begin{equation}\label{eq:echo-default}
 d_h=\sum_{j\geq1}\omega\bigl(1-j(1-b_h)\bigr).
\end{equation}
Here terms outside the finite set just described are zero.  If \(b_h=0\),
then \(d_h=0\).  If \(b_h>0\), its first term gives
\(d_h\geq\omega(b_h)\geq1\).  Each summand is nonincreasing as \(b_h\)
decreases, and no new positive index can appear.  Hence every \(d_h\) is a
nonincreasing nonnegative integer along the chain.  After discarding
finitely many further junctions, all \(d_h\) are constant, including
throughout each intervening finite string of moves.  Indeed each of the
finitely many integers can drop only finitely often, while coefficients
stay unchanged during the small steps.  We henceforth use this truncated
chain and these constant values \(d_h\).

On a fixed model, the generic echo contribution is the only contribution
that can recur over infinitely many surface centers.  We now justify this
claim.  All places in the following proposition are
global divisorial places.

\begin{proposition}[Finiteness below discrepancy two]
\label{prop:difficulty-finite}
Let \((Y,\Theta+\M)\) be a model of the truncated terminal chain.
Among the places exceptional over \(Y\) with discrepancy less than two,
there are only finitely many whose centers have codimension at least three.
For each fixed irreducible compact analytic surface \(V\subset Y\), there
are only finitely many whose center is \(V\).

There is a finite set \(\mathcal V_Y\) of irreducible compact analytic
surfaces with the following additional property.  If
\(V\notin\mathcal V_Y\), its general point lies in a smooth open where
\(\M\) descends and either no label occurs or exactly one smooth label
\(S_h\) occurs.  There are no positive-weight places centered on \(V\)
in the first case or when \(b_h=0\).  When \(b_h>0\), they are exactly the
echoes \(E_{V,h,j}\) with \(j(1-b_h)<1\).
\end{proposition}

\begin{proof}
Take a projective resolution \(\pi:W\to Y\) with \(W\) smooth and carrying \(\M\),
with divisorial exceptional locus, and resolving all labels.  Further
blowups of intersections in the simple normal crossing divisor make the
strict transforms of the labels smooth and pairwise disjoint, while keeping
the full divisor simple normal crossing.  Write
\[
 K_W+\Theta_W+M_W=\pi^*(K_Y+\Theta+M_Y).
\]
Every \(\pi\)-exceptional component of \(\Theta_W\) has coefficient
\(1-a(E;Y,\Theta+\M)<0\), by generalized terminality.  The only possible
positive components of \(\Theta_W\) are therefore the strict transforms
of the positive labels.  On higher models the nef datum pulls back from \(W\), so the
discrepancy calculation is the ordinary one for the signed boundary
\(\Theta_W\).

Consider a place \(E\) that is also exceptional over \(W\), and let its
center on \(W\) have codimension \(c\geq2\).  At a general smooth point
of this center choose normal parameters \(x_1,\ldots,x_c\).  The Jacobian
formula on a smooth model carrying \(E\) gives
\[
 a(E;W,0)\geq\sum_{i=1}^c\ord_E(x_i),
 \qquad \ord_E(x_i)\geq1.
\]
Indeed, in the pullback of a local volume form at most one differential
factor can lose one order along \(E\); adding one to the Jacobian order
gives the displayed bound.  If the center lies in no positive component of
\(\Theta_W\), the signed boundary can only raise this discrepancy, which
is at least two.  Otherwise it lies in a unique positive smooth label of
coefficient \(b<1\).  Taking \(x_1\) as its local equation and discarding
the negative boundary terms gives
\begin{equation}\label{eq:local-discrepancy-estimate}
 a(E;Y,\Theta+\M)
 \geq (1-b)\ord_E(x_1)+\sum_{i=2}^c\ord_E(x_i).
\end{equation}
A discrepancy less than two thus forces \(c=2\) and containment in a
positive label.  A label of coefficient zero cannot give such a place.

If this codimension-two center is contained in \(\Exc(\pi)\), it is a
component of the intersection of its label with one of the finitely many
exceptional divisors.  There are only finitely many such centers.  If it
meets the isomorphism open, its general point maps isomorphically to a
surface downstairs, and the center is the strict transform of that surface.
Thus a center mapping to codimension at least three belongs to the first
finite list, and for any fixed surface downstairs there are only finitely
many possible centers on \(W\).

There are only finitely many places of discrepancy less than two whose
center on \(W\) equals a fixed one of these codimension-two centers.  To see this, discard the negative
components of \(\Theta_W\); this lowers discrepancies, so it suffices to
work with the single smooth label of coefficient \(b\).  Blow up the center
generically.  The new exceptional coefficient is \(-(1-b)\).  A further
place with discrepancy less than two has its center in that exceptional
divisor; the estimate \eqref{eq:local-discrepancy-estimate} also forces a
codimension-two center in the strict label.  It must therefore follow their
intersection above the general point of the original center.  At stage
\(j\) the most recent exceptional coefficient is \(-j(1-b)\), and the
older exceptional divisors do not meet this generic intersection.  A place
above it that is still exceptional has discrepancy at least
\[
 (1-b)+1+j(1-b)=1+(j+1)(1-b).
\]
This reaches two after finitely many stages.  The same argument shows that,
where no other boundary or nef data enter, the places below two are exactly
the echoes in \eqref{eq:echo-discrepancy}.

This generic computation identifies global places.  At each global blowup
there is a unique irreducible component of its exceptional divisor that
dominates the dense smooth part of the center.  The global analytic
intersection of that prime with the strict label then has a unique
irreducible component dominating the same part.  These global components
have coherent ideals, which define the next global blowups.  Conversely,
one can follow any given global place on common higher models through these
blowups; when its center becomes a divisor on a normal model, that divisor
represents the place.  Thus the computation does not count different local
pieces of one place separately.

Finally include the finitely many divisors on \(W\) that are exceptional
over \(Y\).  The complement of the isomorphism open, and the loci where a
label fails to be smooth and alone or the nef datum fails to descend, are
contained in a proper analytic subset of \(Y\) of codimension at least
two.  Such a subset has only finitely many surface components on the compact
fourfold.  Take these components for \(\mathcal V_Y\), enlarging the set
by the finitely many surface images already encountered if necessary.
Outside them the preceding one-label calculation applies, or there is no
positive label and no place of positive weight.  This proves all assertions.
\end{proof}

For an irreducible compact analytic surface \(V\subset Y\), let
\begin{equation}\label{eq:difficulty-branches}
 r_{V,h}=\#\{D:\ D\text{ is a prime divisor of }S_h^\nu
                         \text{ mapping onto }V\}.
\end{equation}
This number is finite, because normalization is finite and the inverse
image of \(V\) has finitely many components on the compact normalization.
It is zero if \(V\not\subset S_h\).  It counts prime divisors on the
normalization, without counting their possible multiple sheets over \(V\)
as separate divisors.

Every place with center a surface, curve, or point on the normal fourfold
is exceptional over it: a place with a divisor as center is represented by
that divisor.  Hence the two raw sums below involve only exceptional
places.  Proposition~\ref{prop:difficulty-finite} makes the local sum
\(\sum_{\cent_Y(E)=V}w(E)\) finite after zero terms are omitted.  Although
this raw sum is finite at each fixed surface, the same nonzero contribution
\(d_h\) can recur at infinitely many surfaces in a label.  We therefore
subtract the branch contributions at each surface before summing over
surfaces.  Define the integer
\begin{equation}\label{eq:difficulty-bracket}
 q_Y(V)=\sum_{\cent_Y(E)=V}w(E)-\sum_{h\in\mathcal H}r_{V,h}d_h.
\end{equation}
Only finitely many \(q_Y(V)\) are nonzero.  Indeed, outside the finite set
in Proposition~\ref{prop:difficulty-finite}, a surface on a positive label
has raw sum \(d_h\), with \(r_{V,h}=1\) and all other branch counts zero.
A surface generically on a sole zero-coefficient label has both raw sum and
default zero, as does a surface away from all labels.  These cases exhaust
the surfaces outside that finite set.

We may therefore define the difficulty by
\begin{equation}\label{eq:difficulty}
 \begin{aligned}
 \mathfrak D(Y,\Theta+\M)
 ={}&\sum_{\codim_Y\cent_Y(E)\geq3}w(E)
   +\sum_{\substack{V\subset Y\\\dim V=2}}q_Y(V)\\
 &+\sum_{h\in\mathcal H}d_hc_2(S_h^\nu),
 \end{aligned}
\end{equation}
where the surface sum is over irreducible compact analytic surfaces.
The first sum has finitely many nonzero terms by
Proposition~\ref{prop:difficulty-finite}; the second has just been shown
to have finite support; and the last is finite because \(\mathcal H\) is
finite and every analytic cycle rank is finite.  Thus \(\mathfrak D\) is
an integer.  The subtraction in \eqref{eq:difficulty-bracket} is performed
at each surface before the second sum is taken.  Both a local bracket and
the difficulty on an intermediate model are at this point allowed to be
negative.

\subsection{The lower bound at the junctions}

The possible negative part of a surface bracket comes only from rounding
the contributions of varying labels.  For a surface \(V\), put
\[
 r_V^{\mathrm{var}}=\sum_{h\text{ varying}}r_{V,h}.
\]
We first prove, on every current terminal model, that
\begin{equation}\label{eq:difficulty-local-bound}
 q_Y(V)\geq-\max\{r_V^{\mathrm{var}}-1,0\}.
\end{equation}

The underlying fourfold is ordinarily terminal by
Lemma~\ref{lem:ordinary-singularities}, hence smooth at general points of
\(V\).  Blow up \(V\) generically to obtain a global exceptional
place \(E_V\).  For each label let \(m_h\) be its multiplicity transverse
to \(V\) at a general point, taking \(m_h=0\) if it does not contain
\(V\).  The ordinary log discrepancy of \(E_V\) is two.  On a common
carrier \(\pi:W\to Y\) carrying this place, the nef defect
\(J=\pi^*M_Y-M_W\) is effective by Lemma~\ref{lem:nef-defect}.  If \(j_V\) is its coefficient at
\(E_V\), the discrepancy formula and generalized terminality give
\begin{equation}\label{eq:difficulty-first-blowup}
 a(E_V;Y,\Theta+\M)=2-\sum_hm_hb_h-j_V>1,
 \qquad j_V\geq0.
\end{equation}
In particular \(\sum_hm_hb_h<1\).

We also have \(r_{V,h}\leq m_h\).  To see this analytically, work at a
general smooth point of \(V\).  Each prime divisor of \(S_h^\nu\) counted
by \(r_{V,h}\) gives at least one point of the finite normalization over
that point of \(V\); the points supplied by distinct prime divisors can be
taken distinct after avoiding their special loci.  They correspond to
distinct local analytic branches of \(S_h\), each containing the germ of
\(V\).  Each branch contributes at least one to the transverse multiplicity
\(m_h\).  A normalization divisor with several sheets may supply more
than one local branch, which preserves the stated inequality.

Let \(t=r_V^{\mathrm{var}}\), and list the numbers \(Nb_h\) for varying
labels, repeating each one \(r_{V,h}\) times, as
\(x_1,\ldots,x_t\).  The fixed contribution
\(k=N\sum_{h\text{ fixed}}r_{V,h}b_h\) is a nonnegative integer.  For
\(t\geq1\), the ceiling inequality gives
\[
 \left\lceil k+\sum_{\ell=1}^t x_\ell\right\rceil
 \geq k+\sum_{\ell=1}^t\lceil x_\ell\rceil-(t-1).
\]
For \(t=0\) the left side is \(k\).  Using effectivity,
\(m_h\geq r_{V,h}\), and \eqref{eq:difficulty-first-blowup}, we obtain in
both cases
\begin{equation}\label{eq:difficulty-first-weight}
 w(E_V)\geq\sum_h r_{V,h}\omega(b_h)
              -\max\{r_V^{\mathrm{var}}-1,0\}.
\end{equation}
No rounding loss is needed for a fixed label because \(Nb_h\) is integral.

It remains to supply the terms of \(d_h\) with \(j\geq2\).  Such a term
is positive only if \(b_h>1/2\); when \(b_h=1/2\), the second echo has
discrepancy exactly two and weight zero.  The strict inequality
\(\sum_hm_hb_h<1\) implies that at most one label through \(V\) has
coefficient greater than \(1/2\), and its multiplicity is one.  This label
is smooth at general points of \(V\) and has \(r_{V,h}=1\).  Use its
successive generic blowups above \(V\).  For each \(j\geq2\) with a
positive default, the resulting global place is distinct from \(E_V\) and
from all the others.  With this label alone its discrepancy is
\(1+j(1-b_h)\).  The other effective boundary components and the effective
nef defect only lower the discrepancy.  Its actual weight is therefore at
least \(\omega(1-j(1-b_h))\).  These places supply every remaining
default.  Adding their weights to \eqref{eq:difficulty-first-weight}
proves \eqref{eq:difficulty-local-bound}, including the cases with no
varying label or with zero coefficients.

\begin{proposition}[Lower bound at crepant junctions]
\label{prop:difficulty-lower}
At every crepant junction \(h_i:(Y_i,\Theta_i+\M)\to(X_i,B_i+\M)\) of
the truncated chain,
\[
 \mathfrak D(Y_i,\Theta_i+\M)\geq0.
\]
\end{proposition}

\begin{proof}
Write \(Y=Y_i\) and let
\(Q=\bigcup_{h\text{ varying}}S_h\) with its reduced structure.  If
\(Q\) is empty, \eqref{eq:difficulty-local-bound} makes every surface
bracket nonnegative, and the other terms of \eqref{eq:difficulty} are
nonnegative.  Suppose \(Q\ne\varnothing\).  At this junction the map
\(h_i\) is projective bimeromorphic from an ordinarily terminal compact
K\"ahler fourfold, and \(\dim h_i(Q)\leq1\) by
Proposition~\ref{prop:surface-finiteness}.  Thus Theorem~\ref{thm:branch}
applies.  For a surface \(V\subset Q\), its branch count in that theorem
is \(r_V(Q)=r_V^{\mathrm{var}}\geq1\); the last inequality follows from
the finite surjective normalizations.  Surfaces outside \(Q\) have
\(r_V^{\mathrm{var}}=0\).  The total possible deficit in
\eqref{eq:difficulty-local-bound} is consequently \(\beta(Q)\).

Every varying label has \(d_h\geq1\).  The places counted by
\(\tau_Y(Q)\) in Theorem~\ref{thm:branch} also each supply a weight of at
least one in the first sum of \eqref{eq:difficulty}.  In fact such a place
\(E\) has ordinary log discrepancy two and center a singular curve in
\(Q\).  It is exceptional over \(Y\).  Its center is contained in a
varying label of positive coefficient, and that effective label is
\(\Q\)-Cartier on the strong model.  Its order at \(E\) is therefore
strictly positive.  On a common carrier the discrepancy formula reads
\[
 a(E;Y,\Theta+\M)
 =2-\ord_E(\Theta)-\operatorname{coeff}_E(J)<2,
 \qquad J\geq0.
\]
The count \(\tau_Y(Q)\) uses distinct global places.  They occur as
distinct terms of the codimension-at-least-three sum, and none is a place
in a surface bracket.  Theorem~\ref{thm:branch} now gives
\[
 \beta(Q)
 \leq\sum_{h\text{ varying}}c_2(S_h^\nu)+\tau_Y(Q)
 \leq\sum_{h\in\mathcal H}d_hc_2(S_h^\nu)
       +\sum_{\codim_Y\cent_Y(E)\geq3}w(E).
\]
Summing \eqref{eq:difficulty-local-bound} and substituting this inequality
in \eqref{eq:difficulty} proves the assertion.  The branch inequality and
this lower bound have been used only at the junctions.
\end{proof}

\subsection{Exact change along the chain}

For a small step, the change in the normalization cycle ranks will cancel
the change in the surface defaults.  We establish this equality before
comparing the difficulties.

\begin{lemma}[Normalization cycle ranks]\label{lem:difficulty-ranks}
Consider a small step of the terminal chain, with exceptional loci
\(L,L^+\),
\[
 Y\xrightarrow{f}Z\xleftarrow{f^+}Y^+.
\]
For each label \(h\), let \(e_h\) and \(e_h^+\) be the numbers of prime
divisors in \(S_h^\nu\) and \((S_h^+)^\nu\) lying over \(L\) and
\(L^+\), respectively.  Then these numbers are finite and
\begin{equation}\label{eq:difficulty-rank-change}
 \begin{gathered}
 e_h=\sum_{\substack{V\subset L\\\dim V=2}}r_{V,h},\qquad
 e_h^+=\sum_{\substack{V^+\subset L^+\\\dim V^+=2}}r_{V^+,h}^+,\\
 c_2(S_h^\nu)-c_2((S_h^+)^\nu)=e_h-e_h^+.
 \end{gathered}
\end{equation}
The surface sums run over irreducible compact analytic surfaces.
\end{lemma}

\begin{proof}
Let \(T_h\) be the normalization of the common image of the two labels in
\(Z\).  Restricting the projective maps, composing with finite
normalization, and factoring through \(T_h\) gives projective
bimeromorphic maps of normal compact threefolds
\[
 S_h^\nu\longrightarrow T_h\longleftarrow (S_h^+)^\nu.
\]
They identify over the common isomorphism open.  By
Lemma~\ref{lem:small-comparison}, the common exceptional image
\(A=f(L)=f^+(L^+)\) has dimension at most one.  A normalization divisor
lying over \(L\) has a surface as its image in \(S_h\), because
normalization is finite, and its image in \(T_h\) lies over \(A\).
It is therefore contracted by the displayed map.  Conversely, the map is
an isomorphism off the preimage of \(A\), so every contracted prime divisor
lies over \(L\).  These are exactly the primes counted in the first sum
in \eqref{eq:difficulty-rank-change}.  The same reasoning holds on the
positive side.  The exceptional loci have finitely many surface components,
which also proves finiteness.

We use the following independence fact.  For a projective bimeromorphic
map \(g:S\to T\) of normal compact threefolds, the classes in
\(H_4(S,\Q)\) of the prime divisors contracted by \(g\) are linearly
independent.  Take a resolution \(r:\widetilde S\to S\) that is projective,
has smooth source, and is an isomorphism over \(U=S_{\mathrm{reg}}\).  Put
\(R=r^{-1}(S\setminus U)\).  Normality gives \(\dim(S\setminus U)\leq1\),
and \(\dim R\leq2\).  If a rational combination of the contracted divisor
classes vanished on \(S\), the corresponding combination of strict
transforms would restrict to zero in \(H_4^{\BM}(U,\Q)\).  The localization
sequence
\[
 H_4(R,\Q)\longrightarrow H_4(\widetilde S,\Q)
                  \longrightarrow H_4^{\BM}(U,\Q)
\]
then shows that its class is a rational combination of the classes of the
surface components of \(R\); those classes span \(H_4(R,\Q)\).  Subtract
that combination.  We obtain a homologically trivial rational Cartier
divisor on \(\widetilde S\), all of whose components are exceptional over
\(T\).  For the components of \(R\), this follows since their images in
\(S\), hence in \(T\), have dimension at most one.  On the smooth compact
resolution, homological triviality makes this divisor numerically trivial
on every curve by the ordinary class pairing.  Analytic negativity in both
signs, in the form of \cite[Lemma~2.2]{CanonicalKahler}, makes the
divisor zero.  The strict transforms of the original divisors are distinct
from the resolution-exceptional ones, so all original coefficients vanish.
This proves independence without any factoriality assumption on \(S\).

For either label normalization, restrict its span of compact analytic
surface classes to the common open over \(T_h\setminus\nu^{-1}(A)\),
where \(\nu:T_h\to f(S_h)\) is normalization.  The complement of this
open in either normalization has dimension at most two.  Localization shows that the kernel of this
restriction is spanned by its surface components, precisely the contracted
prime divisors just counted.  Their independence makes the kernel dimension
\(e_h\) on one side and \(e_h^+\) on the other.  The two images agree:
any compact analytic prime surface not removed has a strict transform on
the other normalization through their proper graph over \(T_h\), and the
two classes restrict to the same class on the common open.  The same holds
in the reverse direction.  Subtracting the dimensions of these two cycle
spans proves the last equality in \eqref{eq:difficulty-rank-change}.
\end{proof}

\begin{proposition}[Change of the difficulty]\label{prop:difficulty-change}
The difficulty is nonincreasing under every coefficient change and every
small step of the truncated terminal chain.  More precisely:
\begin{enumerate}
\item Suppose \(\Theta'\leq\Theta\) is a coefficient change on the same
model \(Y\), and write \(w(E)\) and \(w'(E)\) for its old and new weights.
The set
\[
 \mathcal C=\{E:\ E\text{ exceptional over }Y,\ w(E)\ne w'(E)\}
\]
is finite, and
\begin{equation}\label{eq:difficulty-coefficient-change}
 \mathfrak D(Y,\Theta+\M)-\mathfrak D(Y,\Theta'+\M)
 =\sum_{E\in\mathcal C}\bigl(w(E)-w'(E)\bigr)\geq0.
\end{equation}
\item For a small step \(Y\to Z\leftarrow Y^+\) with exceptional loci
\(L,L^+\), write \(w(E),w^+(E)\) for the two weights.  The set
\[
 \mathcal P=\{E:\ E\text{ exceptional over }Y,\
       \cent_Y(E)\subset L,\ w(E)+w^+(E)>0\}
\]
is finite.  Its center condition is equivalently
\(\cent_{Y^+}(E)\subset L^+\), and
\begin{equation}\label{eq:difficulty-small-change}
 \mathfrak D(Y,\Theta+\M)-\mathfrak D(Y^+,\Theta^++\M)
 =\sum_{E\in\mathcal P}\bigl(w(E)-w^+(E)\bigr)\geq0.
\end{equation}
\end{enumerate}
In either move, a strict decrease in the weight of any exceptional place
forces a strict decrease in the difficulty.
\end{proposition}

\begin{proof}
For a coefficient change, the removed boundary \(\Theta-\Theta'\) is
effective and \(\Q\)-Cartier on the strong model.  For every place,
\[
 a(E;Y,\Theta'+\M)-a(E;Y,\Theta+\M)
 =\ord_E(\Theta-\Theta')\geq0.
\]
Thus \(w(E)\geq w'(E)\).  The model, centers, branch counts, and
normalization ranks are unchanged, and the defaults \(d_h\) are constant
by the choice of the tail.  At a fixed surface the two raw sums are finite,
and therefore
\[
 q_Y(V)-q'_Y(V)=\sum_{\cent_Y(E)=V}\bigl(w(E)-w'(E)\bigr).
\]
Let \(\Sigma\) be the finite union of the supports of the two functions
\(q_Y\) and \(q'_Y\) on surfaces.  Outside \(\Sigma\) the displayed
sum is zero.  Every summand is nonnegative, so no exceptional place with
such a surface center can lose weight.  At each surface of \(\Sigma\)
only finitely many places have positive old or new weight, and there are
only finitely many positive weights at centers of codimension at least
three.  This proves that \(\mathcal C\) is finite.  Summing the displayed
identity over \(\Sigma\), adding the finite comparison in codimension at
least three, and canceling the unchanged rank term gives
\eqref{eq:difficulty-coefficient-change}.

For a small step, Lemma~\ref{lem:small-comparison} says that exceptionality
of a global place is the same on the two sides, that the two center
containment conditions in the statement are equivalent, and that
\(w(E)\geq w^+(E)\).  If a center is not contained in the exceptional
locus, its general point corresponds on the common isomorphism open and its
discrepancy is unchanged.  The set \(\mathcal P\) is finite: each
exceptional locus has dimension at most two and has finitely many surface
components; a surface center contained in it is one of these components.
Proposition~\ref{prop:difficulty-finite} gives finiteness above each such
surface and finiteness of all positive-weight places with lower-dimensional
centers, on both sides.

Surface centers not contained in \(L\) correspond by proper strict
transform to those not contained in \(L^+\).  At their general points the
places, discrepancies, and normalization primes over the surfaces
correspond, so the two local brackets agree.  The analogous correspondence
holds for the raw terms at lower-dimensional centers outside the exceptional
loci.  In the exceptional loci, every contributing place is counted exactly
once in the raw part on each side of \eqref{eq:difficulty}.  This remains
true if its center changes from a surface to a curve or point, or conversely.
The finite comparison of those raw parts is therefore
\(\sum_{E\in\mathcal P}(w(E)-w^+(E))\).  The subtracted defaults on the
exceptional surfaces have difference \(-\sum_h d_h(e_h-e_h^+)\).  Hence
\begin{align*}
 &\mathfrak D(Y,\Theta+\M)-\mathfrak D(Y^+,\Theta^++\M)\\
 &\quad=\sum_{E\in\mathcal P}\bigl(w(E)-w^+(E)\bigr)
       -\sum_h d_h(e_h-e_h^+)
       +\sum_h d_h\bigl(c_2(S_h^\nu)-c_2((S_h^+)^\nu)\bigr)\\
 &\quad=\sum_{E\in\mathcal P}\bigl(w(E)-w^+(E)\bigr),
\end{align*}
where the last equality is Lemma~\ref{lem:difficulty-ranks}.  This is
\eqref{eq:difficulty-small-change}.  Every summand is nonnegative.  A
strictly decreasing exceptional weight must have its center in the
exceptional loci, since weights agree off them, and is in \(\mathcal P\).
A positive weight that becomes zero is included by the definition of
\(\mathcal P\).  The same positivity is immediate in
\eqref{eq:difficulty-coefficient-change}; this proves strictness in both
cases.
\end{proof}
